% !PRINTED 476
\weberpage{476}
The two generating permutations $\tau, \omega$ belong to the first kind (Bd. I, \S.\,153), and consequently $P_{168}$ is a divisor of the alternating \gesperrt{permutation group} of seven figures.

\begin{center}
\S.\,128.

\gesperrt{Equations of the seventh degree with a group of degree 168.}
\end{center}

We now turn to the general theory of the special kind of equations of degree 7 whose \gesperrt{Galois} group reduces to $P_{168}$.

Let $\lambda_0, \lambda_1, \lambda_2, \lambda_3, \lambda_4, \lambda_5, \lambda_6$ first be arbitrary quantities, and let

\weq{1}{\tau = (\lambda_0, \lambda_1, \lambda_2, \lambda_3, \lambda_4, \lambda_5, \lambda_6)}

be a cyclic permutation of degree 7. If we further introduce the transposition pair

\weq{2}{\omega = (\lambda_1, \lambda_2)\,(\lambda_3, \lambda_5)}

then these two permutations generate, through their compositions and repetitions, the whole group $P_{168}$.

If, following \S.\,72, (17), one composes $\chi$ and $\Theta$ from them, one obtains (since, as already remarked above, the composition of the permutations must take place in the reverse order to that of the substitutions) the following permutations of the seven indices:

\[
\begin{aligned}
\Theta^3 &= \tau^6 \omega \tau \omega = (1, 2, 5, 3)\,(4, 6),
\qquad \Theta = (1, 3, 5, 2)\,(4, 6) \\
\chi &= \omega \tau^3 \Theta \tau^2 = (1, 4, 2)\,(3, 5, 6),
\end{aligned}
\]

expressions which can also be derived very easily from \S.\,127, (7). One sees that $\lambda_0$ remains unchanged by $\omega, \Theta, \chi$, and consequently by the whole octahedral group $P_{24}$ contained in $P_{168}$. It is easy to form a function of the seven quantities $\lambda$ that belongs to the group $P_{168}$.

Take, for example, the product $\lambda_0 \lambda_4 \lambda_6$, which consists of the $\lambda$ left untouched by $\omega$, and form the sum of the products that arise from it by application of the cyclic permutation $\tau$ and its repetitions:

\weq{3}{v = \lambda_0 \lambda_4 \lambda_6 + \lambda_1 \lambda_5 \lambda_0 + \lambda_2 \lambda_6 \lambda_1 + \lambda_3 \lambda_0 \lambda_2 + \lambda_4 \lambda_1 \lambda_3 + \lambda_5 \lambda_2 \lambda_4 + \lambda_6 \lambda_3 \lambda_5.}
